\summarychapter{4}{If the Tiber Floods}{The Persecution of Christians}

\begin{inbrief}
The early church is remembered as an army of martyrs, yet before AD 200 Rome's attempts to silence
Christians were halfhearted and few emperors were bloodthirsty. The chapter explains why a tolerant
empire nevertheless treated Christians as outlaws: their different way of life, the slanders told
about them, the charge of atheism and, above all, their refusal to call Caesar Lord.
\question{Why did Rome persecute Christians, and why is this period remembered as the age of martyrs?}
\end{inbrief}

The chapter opens with Polycarp, the aged bishop of Smyrna. The governor did not want to kill him and
urged him simply to swear by Caesar. Polycarp answered, ``I am a Christian,'' scorned the threat of
beasts and of a fire ``that burns for an hour,'' and was burned at the stake while the crowd shouted
against ``the father of the Christians, the destroyer of our gods.''

\ssection{Rome's Policy}
Rome was tolerant of the religions of conquered peoples. If they added homage to the emperor to their
own rites, it almost never interfered, and for the Jews, with their fierce loyalty to one God, it even
dropped that requirement. While Christians passed for a Jewish sect they shared this immunity. When the
Jews disowned them, the situation changed, and Rome faced a harder case than Judaism: the Jews were a
closed people who did little proselytizing, but the Christians talked constantly of Jesus, meant to win
the whole empire and were spreading fast. They refused to worship the emperor and urged everyone else
to refuse too. From time to time, therefore, they felt the anger of the state and of the people.

\ssection{Reasons for Persecution}
The main cause of hatred was the Christians' distinctive way of life. The New Testament calls them
\emph{hagios}, ``saints'' or holy ones, a word whose root means \emph{different}. People distrust those
who are different, and without any preaching the Christian ethic was a standing criticism of pagan
life. At its center was the rejection of the gods, which made Christians ``enemies of the human race''
and social misfits. Every pagan meal began with an offering; feasts were held in temple precincts ``at
the table'' of a god; a Christian who declined seemed rude. Gladiator shows they rejected as inhuman,
though their pull was strong: Augustine tells how his friend Alypius went with his eyes shut and ended
by yelling with the rest. Fear of idolatry also cost them work --- the mason on a temple wall, the
tailor of priestly robes, the incense maker; Tertullian even forbade schoolteaching because the
textbooks told the stories of the gods. Hospitals were under the god Aesculapius.

Cut off in this way from much of social and economic life, Christians displayed a new attitude to human
life. In a society where a slave could be discarded or killed, Christian masters treated slaves kindly
and gave them equal rights in the church; a former slave, Callistus, became bishop of Rome. Christians
would not expose weak or unwanted infants, and a Christian mother would refuse a pagan father's order
to throw out a baby girl. Against the sexual excess of a decaying society, Paul's teaching that the
body is the temple of the Holy Spirit brought a strict condemnation of unchastity and a sacred view of
family life.

This general hatred explains the first Roman persecution. In 64 a fire burned for six days and laid
most of Rome in ashes, and rumor blamed the emperor Nero. He turned the blame on the Christians, who
were crucified, sewn into animal skins and thrown to dogs, dragged by bulls and burned as torches in
his gardens. Peter and Paul probably died then: Peter, by report, crucified head down at his own
request, Paul beheaded as a Roman citizen. Such outbursts were rare in the first two centuries, and
there were long periods of peace; but persecution hung over the church like the sword of Damocles. An
informer, a riot or a strict governor could bring it down, because the Christian as a Christian was
legally an outlaw. As Tertullian said, public hatred asked only for the confession of the name.

\ssection{Sex and Slander}
A second cause was slander. Secrecy breeds distrust, and since outsiders were barred from Christian
worship, rumor filled the gap. The ``Love Feast'' (Agape) and the holy kiss of peace gave rise to
charges of orgies --- the kiss was eventually almost abandoned --- and the words of the Lord's Supper
about body and blood gave rise to charges of cannibalism. The mob feared that such wickedness would
bring the anger of the gods on everyone who tolerated it.

This gave fair-minded officials a problem. About 112 Pliny, governor in Asia Minor, asked the emperor
Trajan whether the name of Christian alone was punishable. He had been sending citizens who confessed
to Rome and executing the others, sure that such stubbornness deserved it; those who denied Christ and
sacrificed went free. Trajan broadly approved: Christians were too many to be hunted, but when a case
came before an official the crime of being a Christian had to be dealt with. Pliny at least knew the
charge was not immorality or cannibalism. The letters explain why persecution was sometimes sporadic
and sometimes savage.

A third cause was the charge of atheism. Most people could not understand worship without images, and
the gods were seen as patrons of the city whom every loyal citizen honored at civic festivals.
Christians, who would not, seemed disloyal and dangerous, since neglected gods sent disasters.
Tertullian reports the reflex: if the Tiber floods or the Nile fails to rise, if there is drought,
earthquake, famine or plague, the cry goes up, ``Christians to the lion.''

\ssection{Caesar Is Lord}
The supreme cause was emperor worship, which grew slowly. Roman rule brought impartial justice, safe
roads and seas, the \emph{pax Romana}; gratitude turned the spirit of Rome into the goddess Roma, with
temples in Asia Minor by the second century BC, and then saw that spirit embodied in the emperor. The
first temple to an emperor was built at Pergamum in 29 BC. Early emperors hesitated --- Claudius
refused temples --- but officials came to see in Caesar worship the answer to the empire's great
problem, how to unite peoples from the Euphrates to the Irish Sea. No local faith could be universal;
Rome was. Caesar worship became the keystone of imperial policy, organized in every province, and any
rival allegiance came to look like a threat of chaos.

The emperor Decius (249--251) made it universal and compulsory for all except the Jews. Once a year
each citizen had to burn a pinch of incense, say ``Caesar is Lord,'' and receive a certificate; then
he could worship any god he liked. It was a test of political loyalty, and whoever refused was branded
a traitor. Christians prayed for the emperor but would not pray to him. The praises on Roman coins for
each new emperor --- Lord of the earth, invincible, worthy to inherit the kingdom --- were the very
language Christians used of Christ, so the question was who was worthy to rule the universe, Caesar or
Christ. The one thing no Christian would say was ``Caesar is Lord.'' To Romans this was insane
stubbornness and self-confessed disloyalty; a pinch of incense would have bought freedom to worship
Christ, and the Christians would not compromise. Rome therefore saw them as potential revolutionaries.
In one sense, Shelley adds, Rome was right, for many Christians saw a cosmic struggle. The Revelation
of John, reflecting the imperial cult in Asia Minor at the end of the first century, traces the
oppression to the devil, the great red dragon, working through two beasts, the imperial power and the
imperial worship. The Christians' only defense was this: they conquered ``by the blood of the Lamb and
by the word of their testimony.''

\begin{keyfacts}
\keypeople{Polycarp of Smyrna; Nero; Peter and Paul; Tertullian; Callistus; Alypius; Pliny and Trajan;
Claudius; Decius.}
\keydates{%
  \kd{29 BC}{First temple to the divinity of the emperor, at Pergamum}
  \kd{AD 64}{Fire of Rome; Nero's persecution; probable martyrdom of Peter and Paul}
  \kd{c. 112}{Pliny asks Trajan how to treat Christians}
  \kd{249--251}{Decius makes Caesar worship compulsory throughout the empire}}
\keyterms{\emph{Hagios} (holy, different); Agape (Love Feast); kiss of peace; atheism (worship without
images); \emph{pax Romana}; goddess Roma; emperor worship; ``Caesar is Lord.''}
\end{keyfacts}
